% from main.tex
%%%%%%%%%%  scan idx 112  |  folio 99  %%%%%%%%%%
\[
\begin{tikzcd}
& G \arrow[dl,"+1"'] & \\
E \arrow[rr] && F \arrow[ul]
\end{tikzcd}
\tag{*}
\]
\emph{un
triangle distingué de \(\Dpx(\mathcal C_S)\). Si \(E\) et \(G\) sont \(n\)-pseudo-cohérents (resp. pseudo-cohérents), alors \(F\) est \(n\)-pseudo-cohérent (resp. pseudo-cohérent).}

\emph{Preuve.} a) Trivial à partir des définitions.

b) Il suffit bien entendu de démontrer l'assertion non respée. Le triangle \((*)\) peut encore s'écrire
\[
\begin{tikzcd}
& F \arrow[dl,"+1"'] & \\
G[-1] \arrow[rr] && E \arrow[ul] & .
\end{tikzcd}
\]
D'après a), \(G[-1]\) est \((n+1)\)-pseudo-cohérent. On est donc ramené à prouver que, si \(E\) est \((n+1)\)-pseudo-cohérent et \(F\) \(n\)-pseudo-cohérent, alors \(G\) est \(n\)-pseudo-cohérent. Quitte à se localiser, on peut supposer \(E\) (resp. \(F\)) strictement \((n+1)\)-pseudo-cohérent (resp. \(n\)-pseudo-cohérent). La flèche \(E\xrightarrow{u}F\) de \(\Dpx^-(\mathcal C_S)\) est définie par un couple de flèches de \(\Cpx^-(\mathcal C_S)\)
\[
\begin{tikzcd}
& E' \arrow[dl,"s"'] \arrow[dr,"u'"]\\
E && F ,
\end{tikzcd}
\]
où \(s\) est un quasi-isomorphisme. D'après (2.2 a)), il existe, quitte à localiser, un quasi-isomorphisme \(t:E''\to E'\), où \(E''\) est strictement \((n+1)\)-pseudo-cohérent. On est donc réduit au cas où \(u\) est définie par une «
